\documentclass{article}
\usepackage[utf8]{inputenc}
\usepackage[T1]{fontenc}
\usepackage{geometry}
\usepackage{graphicx}
\usepackage{float}
\usepackage{hyperref}
\usepackage[english,indonesian]{babel}

\title{Optimasi Pemrosesan Template LaTeX Menggunakan Pendekatan Berbasis AI}
\author{Academic Support Agent}
\date{\today}

\begin{document}
\maketitle

\begin{abstract}

\end{abstract}

\chapter{PENDAHULUAN}


\chapter{TINJAUAN PUSTAKA}


% Original content (commented out):
% \documentclass{article}
\usepackage[utf8]{inputenc}
\usepackage[fleqn]{amsmath}
\usepackage{gensymb}
\usepackage{amssymb}
\usepackage{mathtools}
\usepackage{array}
\usepackage[T1]{fontenc}
\usepackage[margin=1in]{geometry}
\usepackage{lmodern}
\usepackage{tabularx}
\usepackage{pesmath}
\usepackage{fancyhdr}
\usepackage{graphicx}
\graphicspath{{images/}}
 
\pagestyle{fancy}
\fancyhf{}
\rhead{DueDate\\Math 9}
\chead{Assignment}
\lhead{Phillip E. Simons}
\rfoot{Page \thepage}
 

% \cong
% \seg{}
% \angle
% \triangle
% \ray{}
% \circle{}
% \ncong
% \para
% \npara
% \perp
% \nperp
% \congsegs{}{}
% \congangles{}{}
% \congtriangles{}{}
% \ncongsegs{}{}
% \ncongangles{}{}
% \ncongtriangles{}{}
% \parasegs{}{}
% \pararays{}{}
% \paralines{}{}
% \nparasegs{}{}
% \npararays{}{}
% \nparalines{}{}
% \nparasegs{}{}
% \npararays{}{}
% \nparalines{}{}
% \perpsegs{}{}
% \perprays{}{}
% \perplines{}{}
% \nperpsegs{}{}
% \nperprays{}{}
% \nperplines{}{}


% \begin{tabularx}{\linewidth}{l p{0.4\linewidth} | l X l X}
%   \multicolumn{2}{c |}{Statements}&\multicolumn{2}{c}{Reasons}\\
%   \hline
  
%   \hline
%   \multicolumn{2}{c |}{}&\multicolumn{2}{c}{QDE}
% \end{tabularx}



\begin{document}
\begin{flushleft}
1.

\end{flushleft}
\end{document}


\end{document}